\documentclass[12pt]{article}
\usepackage{amsmath,amssymb,amsthm}
\usepackage{mathtools}
\usepackage{geometry}
\geometry{margin=1in}

% Theorem environments
\newtheorem{definition}{Definition}[section]
\newtheorem{lemma}[definition]{Lemma}
\newtheorem{theorem}[definition]{Theorem}
\newtheorem{corollary}[definition]{Corollary}
\newtheorem{proposition}[definition]{Proposition}
\newtheorem{remark}[definition]{Remark}

% Custom commands
\newcommand{\B}{\mathbb{B}}
\newcommand{\C}{\mathbb{C}}
\newcommand{\R}{\mathbb{R}}
\renewcommand{\H}{\mathbb{H}}
\newcommand{\Tr}{\mathrm{Tr}}
\newcommand{\re}{\mathrm{Re}}
\newcommand{\im}{\mathrm{Im}}

\title{Step 2: Biquaternionic Stress-Energy Tensor $\mathcal{T}_{\mu\nu}$ in UBT}
\author{UBT Theory Development}
\date{2026}

\begin{document}

\maketitle

\begin{abstract}
We lift the real stress-energy tensor $T_{\mu\nu}$ computed in Step~1
to the full biquaternionic algebra $\mathbb{B} = \mathbb{C}\otimes\mathbb{H}$.
The biquaternionic stress-energy tensor $\mathcal{T}_{\mu\nu}$ is defined by the
same Hilbert prescription applied to the complex-valued UBT action.
We show its decomposition into kinetic, potential, and gauge pieces,
take the real projection $T_{\mu\nu} = \mathrm{Re}(\mathcal{T}_{\mu\nu})$,
and verify the QED limit $\Theta \to \psi_{\mathrm{Dirac}}\otimes A_\mu$.
\end{abstract}

\tableofcontents

% ------------------------------------------------------------
\section{From Real to Biquaternionic Action}

\subsection{Recap from Step 1}

Step~1 (\texttt{step1\_metric\_variation.tex}) derived, in the real $(-,+,+,+)$
sector, the stress-energy tensor
\begin{equation}
  T_{\mu\nu}
  = T^{(\mathrm{kin})}_{\mu\nu}
  + T^{(\mathrm{pot})}_{\mu\nu}
  + T^{(\mathrm{gauge})}_{\mu\nu}
\end{equation}
where each piece follows from $T_{\mu\nu} = -(2/\sqrt{-g})\,\delta S/\delta g^{\mu\nu}$.

\subsection{Complex-Valued UBT Action}

The full UBT action $S[\Theta]$ in \texttt{appendix\_AA\_theta\_action.tex}
is defined over the biquaternion algebra $\mathbb{B} = \mathbb{C}\otimes\mathbb{H}$.
The field $\Theta: \mathbb{B}^4\times\mathbb{C} \to \mathbb{B}\otimes S\otimes G$
takes values in a biquaternionic bundle and the action
\begin{equation}
  S[\Theta] = S_{\mathrm{kin}} + S_{\mathrm{pot}} + S_{\mathrm{gauge}}
              + S_{\mathrm{boundary}}
\end{equation}
is generally complex-valued (biquaternion-valued in its variation).

% ------------------------------------------------------------
\section{Definition of $\mathcal{T}_{\mu\nu}$}

\begin{definition}[Biquaternionic stress-energy tensor]
\label{def:biquat_T}
The \emph{biquaternionic stress-energy tensor} is:
\begin{equation}
  \label{eq:calT_def}
  \mathcal{T}_{\mu\nu}
  := -\frac{2}{\sqrt{-g}}\,\frac{\delta S[\Theta]}{\delta g^{\mu\nu}}
  \;\in\; \mathbb{B},
\end{equation}
where the variation is taken with respect to the real spacetime metric
$g^{\mu\nu}$ while $\Theta$ is held fixed on its biquaternion fiber.
\end{definition}

\begin{remark}
  Definition~\ref{def:biquat_T} applies the same Hilbert formula as in
  Definition~2.1 of \texttt{step1\_metric\_variation.tex}, but now the
  action $S[\Theta]$ is $\mathbb{B}$-valued.  The variation
  $\delta S/\delta g^{\mu\nu}$ is therefore a $\mathbb{B}$-valued tensor
  field.  The real projection defined below recovers the physical
  stress-energy tensor.
\end{remark}

% ------------------------------------------------------------
\section{Decomposition of $\mathcal{T}_{\mu\nu}$}

\begin{theorem}[Additive decomposition]
\label{thm:decomp}
$\mathcal{T}_{\mu\nu}$ decomposes into contributions from each term in
the action:
\begin{equation}
  \label{eq:calT_decomp}
  \mathcal{T}_{\mu\nu}
  = \mathcal{T}^{(\mathrm{kin})}_{\mu\nu}
  + \mathcal{T}^{(\mathrm{pot})}_{\mu\nu}
  + \mathcal{T}^{(\mathrm{gauge})}_{\mu\nu},
\end{equation}
where each piece is the Hilbert variation of the corresponding action term.
The biquaternionic expressions are:
\begin{align}
  \mathcal{T}^{(\mathrm{kin})}_{\mu\nu}
    &= \Tr_{\mathbb{B}}\!\bigl[(\nabla_\mu\Theta)^\dagger\nabla_\nu\Theta\bigr]
       -\tfrac{1}{2}g_{\mu\nu}
       \Tr_{\mathbb{B}}\!\bigl[(\nabla_\alpha\Theta)^\dagger\nabla^\alpha\Theta\bigr],
       \label{eq:calT_kin} \\
  \mathcal{T}^{(\mathrm{pot})}_{\mu\nu}
    &= -g_{\mu\nu}\,V(\Theta),
       \label{eq:calT_pot} \\
  \mathcal{T}^{(\mathrm{gauge})}_{\mu\nu}
    &= K\Bigl[
         \Tr_{\mathbb{B}}\!\bigl[F_{\mu\alpha}F_\nu{}^\alpha\bigr]
         -\tfrac{1}{4}g_{\mu\nu}
         \Tr_{\mathbb{B}}\!\bigl[F_{\alpha\beta}F^{\alpha\beta}\bigr]
       \Bigr].
       \label{eq:calT_gauge}
\end{align}
Here $\Tr_{\mathbb{B}}$ denotes the biquaternionic trace (over all internal
indices: spinor, gauge, and biquaternion).
\end{theorem}

\begin{proof}
  The action $S[\Theta]$ is a sum of three terms.  Since
  $\delta(A+B)/\delta g^{\mu\nu} = \delta A/\delta g^{\mu\nu}
  + \delta B/\delta g^{\mu\nu}$, the decomposition~\eqref{eq:calT_decomp}
  follows immediately.  The individual formulas~\eqref{eq:calT_kin}--\eqref{eq:calT_gauge}
  are the direct biquaternionic generalisations of the real formulas
  derived in Propositions~2.1, 3.1, and~4.1 of
  \texttt{step1\_metric\_variation.tex}, with $\Tr$ replaced by
  $\Tr_{\mathbb{B}}$.
\end{proof}

\begin{remark}[Boundary term]
  The boundary action $S_{\mathrm{boundary}}$ depends on the extrinsic curvature
  of $\partial\mathcal{M}$ and does not contribute to $\mathcal{T}_{\mu\nu}$
  in the bulk.
\end{remark}

% ------------------------------------------------------------
\section{Real Projection}

\begin{definition}[Physical stress-energy tensor]
\label{def:real_T}
The \emph{physical (real) stress-energy tensor} is:
\begin{equation}
  \label{eq:T_real}
  T_{\mu\nu} := \re(\mathcal{T}_{\mu\nu}),
\end{equation}
where $\re: \mathbb{B} \to \mathbb{R}$ is the real part of the biquaternion
scalar (the $\mathbf{1}$-component with zero imaginary and quaternionic
parts), applied component-wise to the tensor indices $(\mu,\nu)$.
\end{definition}

\begin{lemma}[Real projection commutes with decomposition]
\label{lem:re_commutes}
\begin{equation}
  T_{\mu\nu}
  = T^{(\mathrm{kin})}_{\mu\nu}
  + T^{(\mathrm{pot})}_{\mu\nu}
  + T^{(\mathrm{gauge})}_{\mu\nu},
\end{equation}
where $T^{(\mathrm{X})}_{\mu\nu} = \re(\mathcal{T}^{(\mathrm{X})}_{\mu\nu})$.
\end{lemma}

\begin{proof}
  Linearity of $\re$: $\re(A+B) = \re(A) + \re(B)$.
  Applied to~\eqref{eq:calT_decomp}.
\end{proof}

\begin{remark}[Imaginary components]
  The imaginary part $\im(\mathcal{T}_{\mu\nu})$ represents phase curvature
  and nonlocal energy configurations that do not couple to classical matter.
  It satisfies its own equations of motion in the biquaternionic sector but
  is projected out in the real-valued (observable) limit.
\end{remark}

% ------------------------------------------------------------
\section{QED Limit: Consistency Check}
\label{sec:qed}

\subsection{The QED Limit}

Set $\Theta = \psi_{\mathrm{Dirac}}\otimes A_\mu$, where
$\psi_{\mathrm{Dirac}}$ is a Dirac spinor and $A_\mu$ is an Abelian gauge
field, with all non-Abelian and biquaternionic degrees of freedom set to zero.

In this limit the UBT action reduces to:
\begin{equation}
  S_{\mathrm{QED}} = \int d^4x\,\sqrt{-g}\Bigl[
    \bar{\psi}\bigl(i\gamma^\mu D_\mu - m\bigr)\psi
    - \tfrac{1}{4}F_{\mu\nu}F^{\mu\nu}
  \Bigr],
\end{equation}
where $D_\mu = \partial_\mu + ieA_\mu$ and $F_{\mu\nu} = \partial_\mu A_\nu
- \partial_\nu A_\mu$.

\subsection{Stress-Energy in the QED Limit}

\begin{proposition}[QED stress-energy from UBT]
\label{prop:qed_limit}
In the QED limit, the UBT stress-energy tensor reproduces:
\begin{align}
  T^{(\mathrm{Dirac})}_{\mu\nu}
    &= \tfrac{i}{4}\bar{\psi}\bigl(\gamma_\mu\overset{\leftrightarrow}{D}_\nu
       + \gamma_\nu\overset{\leftrightarrow}{D}_\mu\bigr)\psi
       - g_{\mu\nu}\bar{\psi}(i\gamma^\alpha D_\alpha - m)\psi,
       \label{eq:T_Dirac} \\
  T^{(\mathrm{EM})}_{\mu\nu}
    &= F_{\mu\alpha}F_\nu{}^\alpha - \tfrac{1}{4}g_{\mu\nu}F_{\alpha\beta}F^{\alpha\beta},
       \label{eq:T_EM}
\end{align}
so that $T_{\mu\nu}\big|_{\mathrm{QED}} = T^{(\mathrm{Dirac})}_{\mu\nu}
+ T^{(\mathrm{EM})}_{\mu\nu}$.
\end{proposition}

\begin{proof}
\textbf{Electromagnetic part.}
Setting $K=1$ and taking $\Tr_{\mathbb{B}} \to 1$ (Abelian, no trace needed),
formula~\eqref{eq:calT_gauge} gives exactly~\eqref{eq:T_EM}.
This is the standard Maxwell stress-energy tensor derived from the
Yang--Mills action in the Abelian limit. $\checkmark$

\textbf{Dirac part.}
For $\Theta = \psi_{\mathrm{Dirac}}$ (spin-$\tfrac{1}{2}$ Dirac field,
mass dimension $\tfrac{3}{2}$), the kinetic term
$\Tr_{\mathbb{B}}[(\nabla_\mu\Theta)^\dagger\nabla_\nu\Theta]$ reduces to
$\bar{\psi}\gamma_{(\mu}\overset{\leftrightarrow}{D}_{\nu)}\psi$
(symmetrised).
Combined with the potential term $V(\psi) = -m\bar{\psi}\psi + m\bar{\psi}\psi
= 0$ at the vacuum, and on-shell (Dirac equation
$i\gamma^\mu D_\mu\psi = m\psi$), formula~\eqref{eq:calT_kin} becomes:
\begin{equation}
  T^{(\mathrm{kin})}_{\mu\nu}\big|_{\psi}
  = \tfrac{i}{4}\bar{\psi}\bigl(\gamma_\mu\overset{\leftrightarrow}{D}_\nu
    + \gamma_\nu\overset{\leftrightarrow}{D}_\mu\bigr)\psi,
\end{equation}
which matches the canonical Dirac stress-energy tensor~\eqref{eq:T_Dirac}
on-shell.~$\checkmark$
\end{proof}

\begin{remark}[QED tracelessness check]
  For $m=0$: $g^{\mu\nu}T^{(\mathrm{EM})}_{\mu\nu} = 0$ (conformal invariance
  of Maxwell theory in $d=4$) and $g^{\mu\nu}T^{(\mathrm{Dirac})}_{\mu\nu}
  = -4m\bar{\psi}\psi$ (zero for massless fermions).
  Both are standard results. $\checkmark$
\end{remark}

% ------------------------------------------------------------
\section{Summary}

\begin{theorem}[Biquaternionic stress-energy tensor: summary]
The complete result is:
\begin{equation}
  \label{eq:calT_final}
  \mathcal{T}_{\mu\nu}
  = \Tr_{\mathbb{B}}\!\bigl[(\nabla_\mu\Theta)^\dagger\nabla_\nu\Theta\bigr]
    -\tfrac{1}{2}g_{\mu\nu}\Tr_{\mathbb{B}}\!\bigl[(\nabla_\alpha\Theta)^\dagger
    \nabla^\alpha\Theta\bigr]
    - g_{\mu\nu}V(\Theta)
    + K\Bigl[\Tr_{\mathbb{B}}\!\bigl[F_{\mu\alpha}F_\nu{}^\alpha\bigr]
      -\tfrac{1}{4}g_{\mu\nu}\Tr_{\mathbb{B}}\!\bigl[F_{\alpha\beta}F^{\alpha\beta}
    \bigr]\Bigr],
\end{equation}
with the physical stress-energy tensor $T_{\mu\nu} = \re(\mathcal{T}_{\mu\nu})$.
The QED limit ($\Theta \to \psi_{\mathrm{Dirac}}\otimes A_\mu$) reproduces the
standard electromagnetic and Dirac stress-energy tensors exactly.
\end{theorem}

\begin{remark}[Connection to \texttt{appendix\_R\_GR\_equivalence.tex}]
  Equation~\eqref{eq:calT_final} provides the explicit derivation of
  $T_{\mu\nu} = \re(\mathcal{T}_{\mu\nu})$ that was stated without proof
  in \texttt{appendix\_R\_GR\_equivalence.tex}.
  The Einstein equations with matter are derived in
  \texttt{step3\_einstein\_with\_matter.tex}.
\end{remark}

\begin{thebibliography}{99}
\bibitem{wald}    Wald, R.~M. (1984). \textit{General Relativity}.
                  University of Chicago Press.
\bibitem{ps}      Peskin, M.~E.\ \& Schroeder, D.~V. (1995).
                  \textit{An Introduction to Quantum Field Theory}.
                  Westview Press.
\bibitem{step1}   UBT Theory Development.
                  \texttt{step1\_metric\_variation.tex}, 2026.
\bibitem{appendix_AA}
                  UBT Theory Development.
                  \texttt{appendix\_AA\_theta\_action.tex}, 2025.
\end{thebibliography}

\end{document}
